\documentclass[10pt]{article}
\usepackage[margin=20mm]{geometry}
\usepackage{amsmath,amssymb,mathtools,fontspec,unicode-math,xurl,hyperref}
\IfFontExistsTF{DejaVu Serif}{\setmainfont{DejaVu Serif}}{\setmainfont{Latin Modern Roman}}
\setmathfont{latinmodern-math.otf}
\setlength{\parindent}{0pt}\setlength{\parskip}{6pt}\sloppy
\title{Initial scientific review: M25-0097
Singular algebraic curves and infinite symplectic staircases}
\author{Reviewing assistant}
\date{26 October 2026}
\begin{document}
\maketitle
Version and disposition. Dusa McDuff and Kyler Siegel, arXiv:2404.14702v3, submitted 15 July 2025, supplied PDF dated 16 July 2025 (75 pages). Associated DOI: 10.1007/s00222-025-01359-4. This report binds to the author version; equality with the journal Version of Record is not established. Two Major proof-interface findings and two Minor scope findings are confirmed. No counterexample to the main theorems as intended is established. Human reviews remain empty and the scientific decision is pending.\par
\section*{Claims, definitions, and positive checks}
The paper links singular rational curves in del Pezzo surfaces to ellipsoid-embedding staircases. Upper bounds use inflation after cusp resolution; lower bounds use index-zero unicuspidal curves. Constructions use Orevkov twists and toric degenerations with \(T\)-singularities.\par
The conventions are material: curves are reduced and irreducible; \((p,1)\) permits a smooth distinguished point; sesquicuspidal curves can have ordinary double points away from the cusp, and weak algebraic variants allow further singularities. The normalization \([\omega]=c_1\) gives line area \(3\) in \(\mathbb{CP}^2\). The embedding function is an infimum. Theorem A uses \(c'>c\), not attainment at the limiting boundary. Stabilized lower bounds require the stated semipositivity input, satisfied in the monotone applications.\par
\subsection*{Twists and outer corners}
Sections 1–2, pp. 1–22, were checked for the definitions, index formula, seed curves, anticanonical boundaries, and recurrence. Blowing up a node preserves the reduced total anticanonical boundary. The chain has the required \(-2\) components and terminal \(-1\) component. The complementary birational isomorphism and positivity against the boundary support the Nakai argument; the parity argument distinguishes the degree-eight del Pezzo surfaces. The specific transformed exceptional-class Chern-number calculation on p. 17 was not independently expanded in homology.\par
For a well-placed cusp the two local intersections are \(p,q\), giving \(c_1(C)=p+q\) and index zero. The twist \[(p,q)\mapsto(Kp-q,p),\qquad K=c_1(M)^2-2,\] has determinant one, preserves coprimality, and preserves positivity for \(p>q\), \(K\geq3\). Its ratios agree with Table 2.4.1. The reflection \(R(x)=K+1/(x-K)\) is an involution on its domain. The boundary seed \(p=q=1\) needs the elementary first-step construction rather than a literal invocation of a proposition assuming \(p>q\). The plane conic supplies the other initial branch. In Lemma 2.4.3 the Hessian condition becomes the square of a mixed quadratic coefficient after a pure coefficient vanishes; the resulting linear conditions and Bézout exclusion were traced. No substantive issue is asserted in this part.\par
\subsection*{Inflation}
In Section 3, Euclidean cusp-resolution weights satisfy \(\sum w_i^2=pq\) and \(\sum w_i=p+q-1\). For the first stage of a multiple Puiseux cusp they are multiplied by \(k\), giving a self-intersection drop \(k^2pq\). Reviewer calculations checked 773 coprime cases, including weights \((3,2,1,1)\) for \((5,3)\); these support the Euclidean argument and are not a proof by finite sampling. Resolving the cusp and smoothing positive nodes gives nonnegative self-intersection, the inflation hypothesis. After contraction the cohomology factor is \(1+sc\), while the ellipsoid parameter is \[\frac{\epsilon+ks}{1+cs}\longrightarrow\frac{k}{c}.\] This explains the strict \(c'>c\) in the upper bound. A special two-Puiseux lemma invoked in the more general case needs the first-stage weight argument. The full local analytic normalization, inflation theorem, and relative contraction theorem were not independently reproved.\par
\subsection*{Toric curves, classification, and stabilization}
Sections 4–5 depend on the distinction between a dual Fano polygon and a polygon with the same fan. Balanced exponent vectors make the product parametrization extend at infinity; coprime boundary data rule out a multiple cover. Boundary contact orders give the indicated exponents. The logarithmic derivative has numerator degree at most the number of punctures minus two. In the three-puncture case the rank-two torus map makes the pair of pants immersed. The cited automatic-transversality and deformation inputs were checked at their interfaces, not fully reconstructed. Contact order \(k\) after blowdowns yields \((kq;kp,kp+1)\).\par
In Section 6, for sufficiently small nonzero \(t\), the selected component \(y=0\), \(w^r=\zeta z^a\), \(\zeta^m=-t\), has the required local \((r,a)\) behavior and avoids the other singular chart. The later algebraic-to-symplectic comparison uses F1 below. The first-Hirzebruch-surface classification imports the cited classification of perfect exceptional classes; it was not independently rebuilt. Its explicit index-zero condition is important.\par
In Section 7 the Fibonacci degree/cusp formulas give exactly one residual node by the genus calculation; nine recurrence stages were checked exactly. The degree-three seed still relies on the cited enumerative count and degeneration arguments. Stabilized folding results and conjectural extensions were not treated as new unproved theorems of this paper.\par
\section*{F1 — Major: the non-Gorenstein smoothing interface}
Locations. Eq. (4.1.3), pp. 32–33; Proposition 4.2.7 and proof, pp. 41–42; its application in Theorem 6.1.7, pp. 63–64. Author TeX lines 1752–1757 and 2244–2256.\par
\subsection*{Missing deformation direction}
The claimed miniversal base is \(\mathbb C^{m-1}\) and the polynomial begins \[xy=z^{rm}+C_{m-2}z^{r(m-2)}+\cdots+C_0.\] For \(r>1\) the \(z^{r(m-1)}\) coefficient cannot generally be removed by translating \(z\), since translation breaks the nontrivial \(\mu_r\) action. Set \(m=1,r=2,a=1\). The singularity \(\frac14(1,1)\) has a nontrivial smoothing \[X_t=\{xy=z^2+t\}/\mu_2,\qquad (x,y,z)\mapsto(-x,-y,-z).\] The central quotient is \(\mathbb C^2/\mu_4\) with scalar action: use \((x,y,z)=(u^2,v^2,uv)\) and the residual action induced by multiplication by \(i\) on \((u,v)\). For \(t\ne0\) the cover is smooth and the involution is free. Indeed, a zero gradient or a fixed point forces \(x=y=z=0\), hence \(t=0\). Flatness follows from the hypersurface family and taking invariants of a finite group in characteristic zero. The index-one cover gives a \(\mathbb Q\)-Gorenstein deformation of index two. Thus an actual parameter exists when the printed base has dimension zero.\par
As a primary cross-check, Hacking–Prokhorov, Smoothable del Pezzo surfaces with quotient singularities, arXiv:0808.1550v2, Remark 2.5 and Eq. (1), pp. 4–5, gives the one-dimensional \(\frac14(1,1)\) case and the full range \[xy=z^{dn}+t_{d-1}z^{(d-1)n}+\cdots+t_0.\] The elementary example supplies an independent check.\par
\subsection*{The total-space assumption excludes the intended fibers}
Proposition 4.2.7 assumes the total space \(\mathcal X\) is a complex manifold. Flatness over a disk makes the central fiber a Cartier hypersurface in a smooth threefold, hence Gorenstein. But \[\operatorname{index}\left(\frac1{mr^2}(1,mra-1)\right)
=\frac{mr^2}{\gcd(mr^2,mra)}=r.\] A central singularity with \(r>1\) cannot occur under that hypothesis. In the example, eliminating \(t\) gives total space \(\mathbb C^3/\{\pm1\}\), singular at the origin. This is not a counterexample to the proposition under its own assumptions: it proves that the proposition does not justify its intended non-Gorenstein applications.\par
Counterevidence and impact. Eq. (4.1.5) has a valid product-form smoothing, and the chart in Section 6.3 has a genuine smoothing parameter. These and the known local geometry suggest repair. Section 2’s independent twist construction is unaffected. Nevertheless Theorem 6.1.7 explicitly uses this comparison, hence the corresponding Theorems D/F route contains a demonstrated invalid core proof interface. This warrants Major without asserting that the conclusions are false.\par
Repair and limits. Restore the full equivariant miniversal family. Formulate and prove the comparison for a normal \(\mathbb Q\)-Gorenstein total space with local quotient models. Transport away from singular loci must still give the claimed boundary identifications, extension, tameness, and symplectic class. Replacing the word manifold alone does not verify those steps. A complete repaired comparison theorem is not claimed here.\par
\section*{F2 — Major: scaling in Theorem E breaks the invoked hypotheses}
Locations. Theorem 5.3.1, pp. 50–51; proof of Theorem E, p. 51; author TeX lines 2786–2789. The theorem requires a dual Fano lattice \(T\)-polygon with coprime affine edge lengths. The proof chooses the least \(s>0\) such that \(sQ\) is a lattice polygon, asserts these hypotheses for \(sQ\), and derives \(k=1\) in triangular cases.\par
Take \[Q=\operatorname{conv}\{(-1,-1),(5,-1),(-1,\tfrac12)\}.\] Its primitive inequalities \(x\geq-1\), \(y\geq-1\), \(x+4y\leq1\) all have support height one, so it is dual Fano. The bottom vertices are Delzant; the top has type \(\frac14(1,1)\), a \(T\)-vertex. It is the monotone plane triangle illustrated in Figure 5.1.1, p. 45. Its area is \(9/2\) and the lengths adjacent to the lower-left smooth vertex are \(6\) and \(3/2\), giving inner-corner data \(p/q=4\), \(y=2/3\), also obtained from Table 2.4.1.\par
The minimum scaling about the displayed origin is \(s=2\): the coordinates force \(s\) integral and the half-integral vertex forces it even. But \(2Q\) has affine edge lengths \(12,3,3\), whose gcd is \(3\), and primitive support heights \(2\). It is neither dual Fano nor primitive in edge lengths. The proof’s formula also gives \[k=s/(qy)=2/(2/3)=3,\] not \(1\). This is a concrete failure of the asserted applicability and normalization.\par
Counterevidence and impact. The intended conclusion survives in this example: a conic has self-intersection \(4=pq\) and index \(6\cdot2-2\cdot4-2\cdot1=2\). Translate the smooth vertex to zero and scale by \(2/3\) to obtain vertices \((0,0),(4,0),(0,1)\) and primitive lengths. This diagnoses a possible normalization repair, but the new triangle is still not dual Fano. The underlying toric-curve construction may extend under a weaker assumption on the fan, but that extension is not the stated theorem. Proposition 5.1.2’s visible embeddings remain unaffected.\par
Repair and limits. State and prove a version for translated/scaled lattice \(T\)-polygons with the appropriate fan and deformation properties. Normalize relative to the smooth vertex and primitive lengths, track the rescaled symplectic class, and check the asserted \(k\) in each case. The invalid application lies on the main inner-corner proof route, hence Major. No nonexistence of inner-corner curves is asserted.\par
\section*{F3 — Minor: an additional exceptional pencil parameter}
Locations. Eq. (6.3.1), p. 59; Proposition 6.3.3, p. 60; global-parameter extensions on pp. 61–63. Author TeX lines 3161–3163 and 3186–3206. Choose \(m=1,r=2,a=3\), positive with \(\gcd(r,a)=1\). The setup’s triangle has vertices \((0,0),(0,5),(4,0)\), a Delzant distinguished vertex, an adjacent singularity \(\frac14(1,5)=\frac14(1,1)\), and the third singularity \(\frac15(1,4)\). Even the stronger \(T\)-triangle condition holds. In \(\mathbb{CP}(1,5,2,3)\) the pencil is \[(1+t)xy=w^2+t z^3.\] The denominator in the printed formula is undefined at \(t=-1\). In the natural projective pencil closure the fiber is \(w^2=z^3\). On the smooth chart \(x=1\) its equation is independent of \(y\) and its gradient vanishes along \(z=w=0\). It is singular along \([1:y:0:0]\), including \([1:0:0:0]\), contrary to the assertion that every \(t\ne0,\infty\) has only the designated singularity at \([0:1:0:0]\).\par
Counterevidence, repair, and scope. Lemma 6.4.2 and the proof of Theorem 6.1.7 explicitly use \(t\) sufficiently close to zero, avoiding \(-1\). Exclude \(-1\) from the generic-fiber assertions, or restrict to a small disk near zero. The same global projective pencil cannot have the asserted description at \(-1\). The demonstrated error is auxiliary global scope; no damage to the local construction is proved, so this finding is Minor.\par
\section*{F4 — Minor: missing index zero in the plane warmup}
Location. Lemma 6.5.1, p. 64, author TeX lines 3387–3392. It states that Theorem 2.1.3 gives all data for rational unicuspidal plane curves with one Puiseux pair, without an index condition.\par
The irreducible rational cubic \[C=\{X^3=Y^2Z\},\qquad [s:t]\mapsto[s^2t:s^3:t^3],\] has gradient \((3X^2,-2YZ,-Y^2)\) and its only singularity is \([0:0:1]\), locally of type \((3,2)\). All literal curve hypotheses hold, but degree \(3\) is absent from the odd Fibonacci degrees \(1,2,5,13,\ldots\) in Theorem 2.1.3. Its index is \(6\cdot3-2\cdot3-2\cdot2=8\).\par
Counterevidence, repair, and scope. Theorem 1.2.2(a) already includes this family. The paragraph before the lemma explicitly intends the index-zero list in Theorem 1.2.2(d); Theorem F and the perfect-class definition impose index zero. Insert that condition in the lemma and in the cited symplectic-classification application. This is a false auxiliary scope statement with a precise local repair, not a counterexample to Theorem F, hence Minor.\par
\section*{Coverage, reproducibility, and remaining limits}
Full scientific text on PDF pp. 1–71 was read. The source-image inspection record lists the pages actually opened, including all finding locations and decisive counterevidence. References pp. 72–75 and other source images were not reviewed. The supplied author TeX was statically parsed in full (153 labels and retained targeted contexts); author programs and author TeX were not executed. Missing computational code is not treated as a defect.\par
Reviewer-owned exact computations checked the quotient-smoothing example, singular pencil fiber, cubic, polygon normalization, Euclidean weights, reflection, inflation factor, and Fibonacci genus identities. The all-hypothesis mathematical checks are given above. These do not independently certify external enumerative geometry, automatic transversality, perfect-class classification, or the full analytic deformation and transport arguments.\par
Repository metadata and a limited correction search were checked; exhaustive absence of corrections is not asserted. Primary model cross-check: Hacking–Prokhorov, arXiv:0808.1550v2, pp. 4–5. Accompanying records bind the actual version, current source, unchanged build dependencies, compiler, PDF, and inspection of every final report page.\par
\end{document}
